% Laget automatisk av elevark/lag_elevark.py fra programmering/python-monte-carlo.html. Endringer her blir overskrevet.
\documentclass[11pt]{article}
\usepackage{elevark}
\setlength{\svarhoyde}{1.4cm}
\begin{document}
\elevarkhode{Fysikk 2 · Tenke og gjøre fysikk}{Usikkerhet med Monte Carlo-simulering}{https://lektorbo.no/programmering/python-monte-carlo.html}{Koden på nett}

Regnereglene for usikkerhet er lette å bruke, men hvor kommer de fra, og når slutter de å virke? I en Monte Carlo-simulering lar vi datamaskinen gjøre det samme forsøket 100 000 ganger med tilfeldige målefeil, og ser hvor mye resultatet sprer seg.

\section*{Oppgaver}

\begin{oppgaveliste}
\oppgave \fase{gjett}{Gjett}Hvilken av de to målingene i steg 1 bidrar mest til usikkerheten i $g$? Gir Monte Carlo-simuleringen omtrent samme svar som den enkle regelen eller som kvadratsummen?
\svar
\oppgave \fase{test}{Kjør}Kjør steg 1 et par ganger. Hvor mye endrer svaret seg? Hva skjer med histogrammet hvis du setter $n = 100$?
\svar
\oppgave \fase{test}{Undersøk}Eleven måler tiden for ti svingninger, og usikkerheten i $T$ blir 0,002~s. Endre dT og kjør igjen. Hvilken måling bør nå forbedres?
\svar
\oppgave \fase{forklar}{Forklar}Kjør steg 2. Hver veiing har usikkerhet 0,3~\%, men massen av vannet har over 20~\%. Hvorfor? Hvorfor gir regelen for differanse et større tall enn simuleringen?
\svar
\oppgave \fase{forklar}{Forklar}Kjør steg 3. Hvorfor er gjennomsnittet av rekkevidden mindre enn 10,19~m, og hvorfor går færre enn halvparten av kastene lenger? Sett usikkerheten i farten til 0 for å se bare virkningen av vinkelen.
\svar
\oppgave \fase{forklar}{Vurder}Simuleringen trekker målefeilene fra en normalfordeling, uavhengig av hverandre. Når er ikke det en god modell? Tenk på systematiske feil.
\svar
\oppgave \fase{utfordring}{Utfordring}Lag en Monte Carlo-simulering for $e/m = 2U/(B^{2}r^{2})$ i forsøket med fintrådrøret\footnote{\url{lektorbo.no/opplegg/e-over-m.html}}, med $U = (250 \pm  2)\,\mathrm{V}$, $I = (1{,}40 \pm  0{,}02)\,\mathrm{A}$ og $d = (10{,}0 \pm  0{,}1)\,\mathrm{cm}$, der $B = 0{,}779$ $mT/A \cdot I$. Hvor stor blir usikkerheten i $e/m$, og hvilken måling bidrar mest?
\svar
\end{oppgaveliste}

\end{document}
